\documentclass[11pt]{article}

\usepackage[a4paper,margin=29mm]{geometry}
\usepackage[T1]{fontenc}
\usepackage{lmodern}
\usepackage{microtype}
\usepackage{amsmath,amssymb,amsthm,mathtools}
\usepackage{booktabs}
\usepackage{enumitem}
\usepackage[hidelinks]{hyperref}

\newtheorem{theorem}{Theorem}[section]
\newtheorem{proposition}[theorem]{Proposition}
\newtheorem{corollary}[theorem]{Corollary}
\newtheorem{lemma}[theorem]{Lemma}
\theoremstyle{definition}
\newtheorem{definition}[theorem]{Definition}
\theoremstyle{remark}
\newtheorem{remark}[theorem]{Remark}

\newcommand{\cT}{\mathcal T}
\newcommand{\Pf}{\operatorname{Pf}}

\title{The golden conference operator and its shadow sisters}
\author{Tavis Rudd}
\date{July 2026}

\begin{document}
\maketitle

\begin{abstract}
Let \(C\) be a marked symmetric conference operator on six axes with
\(C^2=5I\), and let \(C_T\) run through its coherent outer six-family.
We prove that middle exterior power, diagonal commutator, golden
eigenspace compression, adjugation, and centered squaring are five forms
of one propagation mechanism.  The diagonal of
\(*\bigwedge^3C_T\) is the signed Joubert cubic vector \(Z\); the same
cubics satisfy \(\Pf[D_x,C_T]=4Z_T(x)\), and over
\(\mathbb Q(\sqrt5)\) they are determinants of the cross-eigenspace
blocks.  The assembled block adjugates recover the Segre--Igusa polar map,
while the six commutator forms give the unique outer-equivariant
synchronized product of pure-spinor big cells.  On balanced Boolean
inputs, \(C^2=5I\) is equivalent to universal maximum-determinant
\(K_{3,3}\) frustration.  Its six sign syndromes form a distance-six
simplex whose transpose gives the order-ten conference operator and the
real \(\operatorname{ETF}(5,10)\).

The construction is reversible at a precise quotient.  The determinant
sextic and the dimer signs both recover the unoriented conference
two-graph; for the Clebsch code this is exactly the invariant unordered
\(10+10\) support split.  A support-half choice is the additional datum
needed for signed odd shadows.  A genuinely source-free product of six
skew-coordinate cells admits no signed synchronized linear lift; the
switching local system in the marked theorem is therefore essential.
We also transport the construction to
Slater amplitudes, Majorana parity walls, and one- and two-factor Abelian
anomaly equations, while keeping these consequences separate from the
operator theorem.
\end{abstract}

\section{Introduction}

A symmetric conference matrix of order six is a small object with an
exceptionally rigid sign system.  We study what happens when its natural
operations are kept together: middle exterior power, commutator with the
diagonal algebra, splitting over \(\mathbb Q(\sqrt5)\), adjugation, and
centered squaring.  The resulting cubics, polar coordinates, determinantal
resolutions, and fermionic amplitudes are images of one marked operator.

The principal theorem is stated on an abstract marked presentation.  This
separates propagation from provenance: a geometric or coding construction
may supply the marked conference class, but the proof below needs only the
operator and its coherent outer marking.

Theorem~\ref{thm:propagation} is the common proof.  Its three mechanisms
are exterior--Pfaffian identification, golden compression, and the two
kernel lines of the assembled differential.  The first turns triangle
holonomy into the Joubert cubics.  The second turns the same Pfaffian into
a \(3\times3\) determinant and its adjugate factorization.  The third
identifies centered squaring as the conormal direction of the Segre cubic.
No exceptional-group or physical interpretation is used in this proof.

Theorem~\ref{thm:frustration} restricts the operator to balanced Boolean
inputs.  Universal maximum determinant is equivalent to the conference
identity, and the six resulting fingerprints form a simplex.  Transposition
then gives the first higher-order conference shadow by a
classification-free Naimark--Gram argument.

Only after propagation do we address provenance.  The projective
deep-hole locus of the Clebsch code is a full rational conic, which forgets
its parent.  The monomial code class retains an unordered split of the
twenty coordinate triples into two \(A_5\)-orbits.  We prove that this split
is exactly the unoriented regular two-graph required by the operator
theorem.  The final section transports the marked cubics to anomaly and
fermionic instruments.  Exceptional, Clifford, code, and lattice parents
are secondary branches and are not needed to follow the main argument.

\section{The marked propagation theorem}

\begin{definition}
A \emph{fully marked golden presentation} consists of an oriented labelled
six-set \(X\); its outer six-set \(\cT\) with the coherent signed outer
\(S_6\)-action; and a coherently oriented family
\((C_T)_{T\in\cT}\) of symmetric zero-diagonal sign matrices satisfying
\(C_T^2=5I\).  When one cross-golden determinant is named, the presentation
also includes a choice of \(\sqrt5\) and compatible orientations of the two
golden determinant lines.  The last choices are unnecessary for rational
Pfaffians, determinants, and polar coordinates.
\end{definition}

For
\(x\in\mathbb Q^X/\mathbb Q\mathbf1\), put
\[
 D_x=\operatorname{diag}(x),\qquad
 K_T=*\bigwedge\nolimits^3C_T,\qquad
 Z_T(x)=\frac14\sum_{|S|=3}(K_T)_{SS}x_S.
\]

\begin{theorem}[Golden operator propagation]
\label{thm:propagation}
For a fully marked golden presentation:
\begin{enumerate}[label=\textup{(\roman*)},leftmargin=*]
\item the vector \(Z=(Z_T)_{T\in\cT}\) carries the signed outer action and
  satisfies
  \[
    \sum_T Z_T=0,\qquad \sum_T Z_T^3=0;
  \]
\item
  \[
    \Pf[D_x,C_T]=4Z_T(x),\qquad
    \det[D_x,C_T]=16Z_T(x)^2;
  \]
\item over \(\mathbb Q(\sqrt5)\), the cross-golden block
  \(B_T(x)=P_{T,-}D_xP_{T,+}\) satisfies
  \[
    Z_T(x)=\pm10\sqrt5\det B_T(x),
  \]
  and its adjugate gives the two conjugate determinantal resolutions and
  rank-one maximal Cohen--Macaulay objects;
\item the centered square
  \[
    W_T=Z_T^2-\frac16\sum_U Z_U^2
  \]
  is the Segre--Igusa polar vector.  If
  \(\widehat W_T=6W_T\),
  \(q_i=6x_i^2-\sum_jx_j^2\), and \(\mathsf A=dZ\) is written on the
  two augmentation quotients in the bases fixed below, then
  \[
    \operatorname{adj}\mathsf A=6\widehat Wq^{\mathsf T}.
  \]
\end{enumerate}
\end{theorem}

The proof first identifies one cubic in three languages.  The polar map is
then forced by the two generic kernel lines of its differential.  A single
exact cofactor fixes the integral normalization in the last identity; the
proof does not infer the identity from numerical interpolation.

\begin{proof}
For a triple \(S=\{i,j,k\}\), expansion of the middle exterior operator
in the oriented triple basis gives
\[
  (K_T)_{SS}=4(C_T)_{ij}(C_T)_{jk}(C_T)_{ki}.
\]
Consequently \(Z_T\) is the triangle-holonomy cubic of the oriented
two-graph defined by \(C_T\).  Coherence of the outer family gives
\(C_{gT}=\operatorname{sgn}(g)M_{g,T}C_TM_{g,T}^{\mathsf T}\), with
\(M_{g,T}\) monomial.  The diagonal algebra is switching-blind, so the
triangle cubics transform by
\[
  gZ_T=\operatorname{sgn}(g)Z_{gT}.
\]
These are the marked Joubert coordinates of six ordered points on
\(\mathbf P^1\) \cite{HMSVOuter,HMSV}; in this normalization their two
relations are \(\sum_TZ_T=\sum_TZ_T^3=0\).  This proves (i), including
the hypotheses under which the classical identification is used.

Put \(A_T=[D_x,C_T]\).  It is alternating, with entries
\((x_i-x_j)(C_T)_{ij}\).  Expanding its Pfaffian over the fifteen perfect
matchings and collecting squarefree cubic coefficients gives
\[
  \Pf A_T=4Z_T.
\]
Squaring the Pfaffian gives \(\det A_T=16Z_T^2\), proving (ii).
This coefficient calculation is also the bridge between the exterior and
commutator descriptions: both read the same triangle holonomies, rather
than merely producing cubics in isomorphic representations.

Now work over \(k(\sqrt5)\).  The projectors
\(P_{T,\pm}=(I\pm C_T/\sqrt5)/2\) split the axis space into two
three-spaces.  Relative to this splitting,
\[
 [D_x,C_T]=
 \begin{pmatrix}
 0&-2\sqrt5\,B_T(x)^{\mathsf T}\\
 2\sqrt5\,B_T(x)&0
 \end{pmatrix}.
\]
Taking determinants yields
\(16Z_T^2=8000\det(B_T)^2\).  The chosen determinant-line orientation
therefore gives \(Z_T=\pm10\sqrt5\det B_T\).  Multiplication by the
scaled adjugate gives
\[
 B_T(\pm10\sqrt5\operatorname{adj}B_T)
 =Z_TI_3.
\]
On \(Z_T=0\), the two kernels of \(B_T\) and \(B_T^{\mathsf T}\) give
the two determinantal small resolutions; conjugating \(\sqrt5\) exchanges
them.  The standard linear resolution of a \(3\times3\) determinant
cokernel gives the conjugate rank-one MCM objects and their rational
rank-two descent.  This proves (iii).

It remains to identify the assembled polar map.  Let \(G_x\) have columns
\(dZ_T|_x\), and pass to the source and target augmentation quotients to
obtain \(\mathsf A(x)\).  The first kernel line comes from congruence.
Writing \(\omega_x=[D_x,C_T]\),
\[
 [D_{x^2},C_T]=D_x\omega_x+\omega_xD_x
 =\left.\frac d{dt}\right|_{0}
 (I+tD_x)\omega_x(I+tD_x).
\]
Pfaffians transform under congruence by the determinant.  On the centered
slice \(\operatorname{tr}D_x=0\), differentiation therefore gives
\(\mathsf A(x)q(x)=0\).  The second kernel line is the differential of
the Segre equation:
\[
 \widehat W(x)^{\mathsf T}\mathsf A(x)
 =6\sum_TZ_T(x)^2dZ_T|_x=0.
\]
A nonzero fourth minor shows that the generic rank is four.  Hence the
adjugate has rank one and is a scalar multiple of
\(\widehat Wq^{\mathsf T}\).  Degree and outer equivariance make the
scalar constant.  In the quotient bases
\(e_i-e_5\) and \(e_T-e_{T_5}\), comparison of one integral cofactor
with the corresponding product gives the scalar \(6\); this exact
normalization check is reproduced in the verification supplement.
Polynomial continuation proves the identity everywhere.

Finally, centered squaring is the projective gradient of
\(\sum_Tz_T^3\) modulo the linear relation \(\sum_Tz_T=0\).  Newton's
identities applied to \(\sum z_T=\sum z_T^3=0\) give
\[
 \sum_TW_T=0,
 \qquad \left(\sum_TW_T^2\right)^2=4\sum_TW_T^4,
\]
the Igusa quartic equation.  Thus the left kernel just found is precisely
the Segre--Igusa polar direction, completing (iv).
\end{proof}

\subsection{The synchronized Cartan cells}

Put \(U=\mathbb Q^X\) and \(A_X=U/\mathbb Q\mathbf1\).  For
\(T\in\cT\) and \(x\in A_X\), the alternating form
\[
 \alpha_T(x)(e_i,e_j)=(x_i-x_j)(C_T)_{ij}
\]
defines the maximal isotropic graph
\[
 L_T(x)=\{u+\iota_u\alpha_T(x):u\in U\}
 \subset U\oplus U^*.
\]
Its Cartan coordinates in the positive half-spin big cell are
\[
 s_T(x)=\sum_{|S|\ \mathrm{even}}
       \Pf(\alpha_T(x)_S)e_S.
\]

\begin{proposition}[Synchronized pure-spinor quotient]
\label{prop:synchronized-spinors}
The map
\[
 \Phi_C:A_X\longrightarrow
 \prod_{T\in\cT}\operatorname{OG}^+(6,12),
 \qquad x\longmapsto(L_T(x))_T,
\]
is outer-\(S_6\)-equivariant.  Its linear big-cell map is unique up to
common scale after the coherent golden marking is fixed.  The six top
coordinates \(p_T=\Pf\alpha_T=4Z_T\) have scheme-theoretic image
\[
 \operatorname{Spec}\frac{\mathbb Q[z_T:T\in\cT]}
 {(\sum_Tz_T,\sum_Tz_T^3)}.
\]
The projectivized top map has reduced base scheme
\[
 \bigcup_{Q\in\binom X4}
 \{[x]:x_i=x_j\text{ for all }i,j\in Q\},
\]
a union of fifteen lines.  The ten strictly semistable \(3+3\)
configurations map to the ten nodes of the Segre cubic.
\end{proposition}

\begin{proof}
The graph of an alternating form is maximal isotropic, and the Cartan map
on this big cell is the principal-Pfaffian expansion
\cite{ChiriviMaffei}.  If \(g\in S_X\), coherence gives a switching
isometry \(M_{g,T}\) with
\[
 C_{gT}=\operatorname{sgn}(g)M_{g,T}C_TM_{g,T}^{\mathsf T}.
\]
The diagonal algebra is switching-blind, so
\[
 \alpha_{gT}(gx)=\operatorname{sgn}(g)
 M_{g,T}\alpha_T(x)M_{g,T}^{\mathsf T}.
\]
This proves equivariance, including the half-spin coordinate signs.

Let
\(
 \mathcal W=\bigoplus_{T\in\cT}\bigwedge^2U_T^*
\)
with this signed monomial action.  On a conjugacy-class representative
\(g\), its character is
\[
 \chi_{\mathcal W}(g)=\operatorname{sgn}(g)
 \sum_{gT=T}\frac{\operatorname{tr}(M_{g,T})^2
 -\operatorname{tr}(M_{g,T}^2)}2.
\]
Its inner product with the augmentation character
\(\#\operatorname{Fix}_X(g)-1\), over the eleven classes of \(S_6\), is
one.  Thus
\[
 \dim\operatorname{Hom}_{S_6}(A_X,\mathcal W)=1.
\]
The displayed map is nonzero and hence spans this Hom-space.  Set
\[
 E=\operatorname{sgn}\otimes\mathbb Q[\cT]_0.
\]
The corresponding cubic multiplicity is also one:
\[
 \dim\operatorname{Hom}_{S_6}(\operatorname{Sym}^3A_X,E)=1.
\]
The top Pfaffians are therefore the unique signed outer cubic covariant;
they are the Joubert coordinates of six ordered points on
\(\mathbf P^1\) \cite{HMSV}.

It remains to identify the whole image ideal.  In the translation gauge
\(x_5=0\), at \((x_0,\ldots,x_4)=(-2,-2,-1,-1,0)\), the target is
\((16,-16,0,0,0,0)\).  A \(4\)-minor of the Jacobian, using target rows
\((0,2,3,4)\) and source columns \((0,1,2,4)\), equals \(-24576\).
Thus the image has affine dimension four.  The prime ideal
\((\sum z_T,\sum z_T^3)\) has the same height as the image ideal.  It is
prime because the projective Segre cubic is irreducible: a reducible cubic
hypersurface would have a positive-dimensional singular intersection,
whereas this cubic has only ten nodes.  Hence the two ideals agree.

Finally, the top map is the equal-weight GIT quotient of six ordered
points on \(\mathbf P^1\) \cite{HMSV}.  Its unstable locus consists of
configurations in which at least four points coincide, giving the fifteen
displayed lines.  Exact radical and minimal-prime computation shows that
this base scheme is reduced.  The \(3+3\) polystable orbits are the ten
nodes.
\end{proof}

The target has another Pfaffian description after a nonstandard
\(S_5\subset S_6\) is chosen: a five-dimensional submodule of
\(\bigwedge^2V^*\), for a six-dimensional representation of a Schur cover,
cuts out the Segre cubic by one Pfaffian equation
\cite{TschinkelZhang}.  That presentation is an equation on the target.
It neither supplies a common Gaussian parent for the six cells above nor
makes the Segre equation a Wick relation.  The quadratic Wick, or
matchgate, identities govern each factor separately
\cite{ChiriviMaffei,CaiGorenstein}; the Segre equation records their golden
synchronization.

The common outer augmentation module also locates the order-ten shadow:
\[
 A_X\xrightarrow{\Phi_C}\prod_T\operatorname{OG}^+(6,12)
 \xrightarrow{\mathrm{top}}E
 \xrightarrow{R^{\mathsf T}/\sqrt{12}}E_{+3}(S_{10}).
\]
The arrows are respectively linear in big-cell coordinates, cubic, and
linear isometric.  They are consecutive functorial operations, not
different coordinate expressions for one linear map.

The uniqueness in Proposition~\ref{prop:synchronized-spinors} is a marked
statement, but it has no hidden deformation parameter.

\begin{corollary}[Rigidity of the marked synchronized lift]
\label{cor:marked-lift-rigidity}
Fix the two six-sets, their outer action, the coherent switching frames, and
the integral axis and skew-coordinate lattices.  Every outer-equivariant
linear map
\[
 A_X\longrightarrow\bigoplus_{T\in\cT}\bigwedge^2U_T^*
\]
whose top Pfaffians are the Joubert cubic vector is a common scalar multiple
of \(x\mapsto([D_x,C_T])_T\).  Primitive integrality leaves only the two
common signs, and the normalization
\(\Pf[D_x,C_T]=4Z_T\) selects one.  Over
\(\mathbb Q(\sqrt5)\), the cross-golden blocks and their adjugates are then
forced; golden conjugation exchanges the two transpose presentations.
\end{corollary}

\begin{proof}
The character calculation in the proof of
Proposition~\ref{prop:synchronized-spinors} makes the relevant Hom-space
one-dimensional.  Primitivity fixes the scalar up to sign, and the displayed
Pfaffian normalization fixes that sign.  The block decomposition in the
proof of Theorem~\ref{thm:propagation} then forces the determinant and
adjugate data.
\end{proof}

This multiplicity-one statement is the order-six member of a classical
matching construction.  For \(2m\) labelled points, the multilinear
\(\operatorname{SL}_2\)-invariants form the Specht module \(S^{(m,m)}\),
generated by bracket monomials modulo the Pl\"ucker relations
\cite{HMSV}.  It first occurs in degree \(m\) in the symmetric algebra of
the axis augmentation module, with multiplicity one: the squarefree
monomials contain \(S^{(m,m)}\) once by Young's rule, while the
nonsquarefree orbit modules are excluded by dominance.  Order six is the
unique nontrivial case in which the equal-weight quotient is a hypersurface,
since
\[
 \operatorname{codim}=\frac1{m+1}\binom{2m}{m}-2m+2
\]
equals one only for \(m=3\).  The Golden input does not create this carrier;
it selects the six-vector simplex of conference functionals on it.

There is also a sharp difference between the intrinsic Pfaffian system and
one chosen cross-golden block.  On \((\mathbf P^1)^6\), with tautological
lines \(L_i\subset V\) and \(E=\bigoplus_iL_i\), the entries
\(C_{T,ij}[ij]\) define an alternating bundle map
\[
 E\longrightarrow E^*\otimes\det V.
\]
Its Pfaffian is therefore intrinsic on the six-point quotient.  After a
seventh point \(p_\infty\) is chosen, it selects an affine coordinate and the
projective cross block \(P_{T,-}D_xP_{T,+}\) is intrinsic on
\(M_{0,7}\to M_{0,6}\).  Under a general projectivity
\(y_i=(ax_i+b)/(cx_i+d)\), however,
\[
 [D_y,C_T]=(ad-bc)\Lambda[D_x,C_T]\Lambda,
 \qquad \Lambda_{ii}=(cx_i+d)^{-1}.
\]
The diagonal matrix \(\Lambda\) preserves a fixed golden eigenspace only
when it is scalar, because every off-diagonal entry of \(C_T\) is nonzero.
Thus forgetting the pole preserves the determinant cubic but not the fixed
golden \(3+3\) splitting or a selected rank-one matrix factorization.  The
continuous pole, support orientation, and golden conjugation are independent
ambiguities.

The collision geometry likewise forms a filtration rather than one master
discriminant.  In the source, pair collisions give fifteen hyperplanes,
the rank-drop scheme of the Joubert Jacobian has twenty triple-collision
planes as its reduced support, and the simultaneous top-Pfaffian base has
fifteen four-collision lines.  The ten \(3+3\) closed orbits map to the ten
Segre nodes.  Appendix~\ref{app:collision-filtration} records the exact
scheme statement and its computational trust boundary.

\subsection{Balanced cuts and the first Naimark shadow}

Let \(B\) be any symmetric zero-diagonal \(6\times6\) sign matrix.  For a
balanced partition \(L\sqcup R=X\), expand \(\det B_{L,R}\) as the six
signed perfect matchings of the crossing \(K_{3,3}\).

\begin{theorem}[Conference rigidity from balanced frustration]
\label{thm:frustration}
The following conditions are equivalent:
\begin{enumerate}[label=\textup{(\roman*)},leftmargin=*]
\item on every balanced cut the six matching terms split \(5{:}1\) by
sign;
\item every cross determinant has absolute value four;
\item \(B^2=5I_6\);
\item after switching so that \(B_{0i}=1\), the negative edges on the
other five vertices form a five-cycle.
\end{enumerate}
The twelve oriented normalized solutions pair under global negation into
six projective fingerprints.  Their ten-cut sign words \(r_T\) satisfy
\[
 \langle r_T,r_U\rangle=
 \begin{cases}10,&T=U,\\-2,&T\ne U,
 \end{cases}
 \qquad
 RR^{\mathsf T}=12\left(I_6-\frac16J_6\right).
\]
\end{theorem}

\begin{proof}
Switch to \(B_{0i}=1\).  Every balanced cut has a representative
\(L=\{0,i,j\}\), so its cross block has first row \((1,1,1)\).  A
\(3\times3\) sign determinant is zero or has absolute value four; it is
nonzero exactly when its six determinant terms have a \(5{:}1\) sign
split.  Thus (i) and (ii) are equivalent.

Fix an internal vertex \(i\).  For every removed vertex \(j\ne i\),
nonvanishing says that the three remaining signs incident with \(i\) are
not constant.  This holds for every \(j\) exactly when two of the four
internal edges at \(i\) are negative.  The negative graph is therefore
two-regular on five vertices, hence a five-cycle.  Conversely, the
one-neighborhoods in a five-cycle give
\((B^2)_{ij}=0\) for \(i\ne j\), while the diagonal entries are five.
The implication (iii) to (iv) already follows from
\((B^2)_{0i}=0\), which forces two signs of each kind at every internal
vertex.  This proves the equivalence.

There are \((5-1)!/2=12\) labelled five-cycles, and complementation pairs
them without fixed points.  On each balanced cut the identity
\(\sum_TZ_T=0\) gives three positive and three negative determinant signs,
so the six row words have zero centroid.  The outer action is
two-transitive on the six sisters, hence all off-diagonal row inner
products are equal.  Since every row has norm squared ten, the common
inner product is \(-10/5=-2\), proving the Gram identity.
\end{proof}

\begin{corollary}[The \(6\to10\) Naimark--Gram lift]
\label{cor:naimark}
The columns of \(R\), modulo sign, are exactly the ten balanced cuts of
the row six-set.  They form the real \(\operatorname{ETF}(5,10)\), and
\[
 S_{10}=\frac{R^{\mathsf T}R-6I_{10}}2
\]
is a symmetric zero-diagonal sign matrix satisfying \(S_{10}^2=9I_{10}\).
Moreover \(R^{\mathsf T}/\sqrt{12}\) maps the six-point augmentation
space isometrically onto the \(+3\)-eigenspace of \(S_{10}\).  Up to row
and column monomial actions, this is the unique sign Gram factorization of
the resulting order-ten switching class.
\end{corollary}

\begin{proof}
The row Gram identity annihilates \(\mathbf1\), so
\(R^{\mathsf T}\mathbf1=0\) and every column is balanced.  Let \(n_C\)
be the multiplicity of the projective cut \(C\), and let \(M\) be the
cut--edge incidence matrix.  Every edge crosses six columns, hence
\(Mn=6\mathbf1\).  A cut has nine edges and two distinct balanced cuts
share five; therefore
\[
 M^{\mathsf T}M=4I_{10}+5J_{10}.
\]
It is nonsingular, so \(n_C=1\) for all ten cuts.  Distinct normalized
cut vectors have inner product of absolute value \(1/3\), and the row
Gram identity gives tightness on the augmentation five-space.  This proves
the ETF statement and shows that \(S_{10}\) is a zero-diagonal sign matrix.

Writing \(G=R^{\mathsf T}R\), one has
\(G^2=R^{\mathsf T}(RR^{\mathsf T})R=12G\).  Hence
\[
 S_{10}^2=\frac{G^2-12G+36I}{4}=9I.
\]
The same calculation identifies the image of
\(R^{\mathsf T}/\sqrt{12}\) with the \(+3\)-eigenspace.  Conversely, any
sign factor of \(6I+2S_{10}\) has row Gram spectrum \(12^5,0\); row
reversal makes its kernel vector \(\mathbf1\), after which the incidence
argument forces the same universal ten-cut matrix.  This proves uniqueness
up to the stated monomial actions.
\end{proof}

The \(36\) extremal balanced cuts of the resulting order-ten operator have
stabilizer \(F_{20}\), the normalizer of a Sylow \(5\)-subgroup \(H<S_6\).
Each \(H\) lies in exactly one standard and one outer \(S_5\): both incidence
counts are \(6\cdot6=36\), the number of Sylow \(5\)-subgroups.  Hence there
is a canonical equivariant bijection
\[
 S_6/F_{20}\simeq X\times\cT.
\]
Projection to \(\cT\) is the unmarked \(36\to6\) return to the Golden
sisters.  It introduces no new projective torsor; orienting the returned
sister still requires the support-half bit.

\section{Recovery and minimal marking}

Let \(X\) be the coordinate six-set of the Clebsch code and let
\(H\cong A_5\) be its monomial permutation group.  The twenty triples in
\(X\) form two complementary \(H\)-orbits
\(\mathcal O_+\) and \(\mathcal O_-\), each of size ten.  The code
canonically supplies the unordered pair
\(\{\mathcal O_+,\mathcal O_-\}\), but does not name a positive half.

\begin{proposition}[The support split is the conference two-graph]
\label{prop:support-two-graph}
Choose one half and put \(t_{ijk}=1\) on it and \(t_{ijk}=-1\) on the
other.  Then there is a symmetric zero-diagonal sign matrix \(C\), unique
up to diagonal switching, such that
\[
 C_{ij}C_{jk}C_{ki}=t_{ijk},\qquad C^2=5I.
\]
Interchanging the two halves replaces the switching class of \(C\) by
that of \(-C\).  Hence the unordered support split determines exactly the
unoriented switching line
\[
 \ell_C=\{[C],[-C]\}.
\]
\end{proposition}

\begin{proof}
The degree-six \(A_5\)-action is two-transitive, hence transitive on pairs
and, by complementation, on four-subsets.  Count incidences between one
ten-element triple orbit and the fifteen four-subsets.  Each triple lies
in three four-subsets, so every four-subset contains two triples from
each orbit.  This is the two-graph parity condition.

Fix \(r\in X\), set \(C_{ri}=1\), and set \(C_{ij}=t_{rij}\) for
\(i,j\ne r\).  The parity identity on \(\{r,i,j,k\}\) gives
\(C_{ij}C_{jk}C_{ki}=t_{ijk}\).  Changing the normalized row changes
only the diagonal switching gauge.  The same incidence count on pairs
shows that two of the four values \(t_{ijk}\), \(k\ne i,j\), have each
sign.  Thus
\[
 (C^2)_{ij}=C_{ij}\sum_{k\ne i,j}t_{ijk}=0
 \quad(i\ne j),
 \qquad (C^2)_{ii}=5.
\]
Finally, \(-C\) negates all triangle holonomies.
\end{proof}

The oriented two-graph has stabilizer \(A_5\).  Its unoriented line has
stabilizer \(N_{S_6}(A_5)\cong S_5\).  Thus the twelve oriented
\(S_6/A_5\)-translates pair under common negation into the outer six-set
\(S_6/S_5\).  After choosing one support half, its \(A_6\)-orbit gives
the coherent oriented six-family.  Different matrix lifts are diagonal
switches, so the covariance in Theorem~\ref{thm:propagation} removes that
gauge.

\begin{theorem}[Recovery--propagation and minimal marking]
\label{thm:recovery-propagation}
A monomial Clebsch code class canonically determines the switching line
\(\ell_C\) and the coherent outer six-family modulo common reversal,
diagonal switching, and permutation.  Consequently:
\begin{enumerate}[label=\textup{(\alph*)},leftmargin=*]
\item the projective Joubert map, Segre point, Segre--Igusa polar map,
  Pfaffian cubic lines, determinant sextics, parity walls, squared Slater
  probabilities, and projective anomaly solution descend to the code
  class;
\item the cross-golden resolutions and rank-one MCM objects descend as an
  unordered conjugate pair, while their rational rank-two descent is
  canonical;
\item choosing one half of the support split is the minimal extra bit for
  a signed odd shadow; choosing \(\sqrt5\) is the independent bit selecting
  one cross-golden summand;
\item the determinant sextic and the relative dimer fingerprint each
  recover \(\ell_C\), so the recovery cycle splits at the unoriented
  two-graph level.
\end{enumerate}
The unlabelled projective deep-hole conic alone does not canonically
determine this input-relative family.  The exact minimal source is the
unordered support two-graph
\[
 (X,\{\mathcal O_+,\mathcal O_-\}).
\]
A geometric Clebsch parent, a monomial code marking, or a support-labelled
decoder is sufficient because each supplies that quotient object.
\end{theorem}

\begin{proof}
Proposition~\ref{prop:support-two-graph} and the \(A_6\)-orbit construction
give the coherent source required by Theorem~\ref{thm:propagation}.
Common reversal negates \(Z\) and every Pfaffian but fixes \(Z^2\), the
centered square, determinants, probabilities, and zero loci.  Golden
conjugation exchanges the two eigenspaces, transposes the cross block, and
exchanges the two resolutions and rank-one sheaves.  These are exactly the
descent claims in (a)--(c).

For (d), unique factorization recovers the cubic line from
\(\det[D_x,C]=16Z_C^2\).  Its triangle coefficients recover \(t\) up to
common negation.  Equivalently, for distinct \(i,j,k,l\),
\[
 [x_i^2x_j^2x_kx_l]\det[D_x,C]
   =32\,t_{ijk}t_{ijl},
\]
so the determinant directly gives the four-cycle holonomies measured by
relative matching signs.  Those holonomies recover the switching line.

For minimality, the bare locus is the full rational conic.  Its
\(\mathrm{PGL}_2(11)\)-action is transitive on the
\([\mathrm{PGL}_2(11):A_5]=22\) Clebsch matching rows, so no equivariant
rule selects a parent.  A factorization-sheet orientation does not repair
this: it selects a decorated row, whereas identifying that row with the
support two-graph is an additional bridge.  Once the unordered support
two-graph is retained, the \(S_5/A_5\cong C_2\) support-half choice is both
sufficient and, for a signed odd shadow, unavoidable.
\end{proof}

The preceding theorem cannot be strengthened by simply deleting its marked
target.  The obstruction can be stated without reference to the Clebsch
source.

\begin{theorem}[Boundary of unmarked reconstruction]
\label{thm:unmarked-boundary}
Let \(X\) be the natural six-set, let \(\cT\) be its outer six-set, put
\(A_X=\mathbb Q^X/\mathbb Q\mathbf1\), and give
\[
 \mathcal W_0=\mathbb Q[\cT]\otimes
              \bigwedge^2(\mathbb Q^X)^*
\]
the source-free product action.  Then
\[
 \dim\operatorname{Hom}_{S_6}(A_X,\mathcal W_0)=1,
 \qquad
 \dim\operatorname{Hom}_{S_6}
 (A_X,\operatorname{sgn}\otimes\mathcal W_0)=0.
\]
The unique ordinary lift repeats the difference form
\(\alpha_{ij}(x)=x_i-x_j\) in all six cells and has zero top Pfaffian.
Thus a nonzero signed Joubert lift requires the coherent switching local
system used in Proposition~\ref{prop:synchronized-spinors}.

The paired natural and outer actions already recover the coherent projective
conference family: for \(T\in\cT\), the derived subgroup of
\(\operatorname{Stab}(T)\cong S_5\) has two complementary ten-element
orbits on \(\binom X3\), whose unordered pair is the regular two-graph
\(\ell_T\).  Hence this action is a viable unmarked input, but it contains
the projective source before the tensor map is evaluated.

Finally, no complex generated by one selected cross-golden block can have
the pair-collision divisor as its first Fitting support.  On the pole-marked
cover that support is the cubic wall \(Z_T=0\); for all six blocks it is
\(\prod_TZ_T=0\), of degree eighteen, rather than the degree-fifteen
pair-collision divisor.  A selected block does not descend to \(M_{0,6}\).
\end{theorem}

\begin{proof}
If \(f_X(g)\) and \(f_{\cT}(g)\) count fixed points in the two six-actions,
then
\[
 \chi_{A_X}(g)=f_X(g)-1,
 \qquad
 \chi_{\wedge^2\mathbb Q^X}(g)
 =\frac{f_X(g)^2-f_X(g^2)}2.
\]
Averaging their product with \(f_{\cT}(g)\), with and without the sign
character, gives \(1\) and \(0\).  The displayed difference form is nonzero,
so it spans the ordinary Hom-space; its matrix is
\(x\mathbf1^{\mathsf T}-\mathbf1x^{\mathsf T}\), of rank at most two.

For the projective recovery, the two triple orbits satisfy the two-graph
parity condition by the four-subset incidence count in
Proposition~\ref{prop:support-two-graph}; conjugation in \(S_6\) gives
coherence.  The Fitting claim follows from
\(\det B_T\doteq Z_T\).  The non-descent is the diagonal-congruence
calculation above: a non-affine \(\Lambda\) does not preserve the golden
eigenspaces.
\end{proof}

The centered-square shadow has one further, geometric information-loss
stratum.  Over an algebraically closed field of characteristic zero, let
\[
 S=\{[z]\in\mathbf P^5:\textstyle\sum_i z_i=\sum_i z_i^3=0\},
 \qquad
 \pi([z])=[\operatorname{center}(z_i^2)].
\]

\begin{proposition}[Fibres of centered squaring]
\label{prop:polar-fibres}
The base locus of \(\pi\) is the set of ten nodes of \(S\).  Let
\(\widetilde\pi:\widetilde S\to I\) be the closure of its graph, equivalently
the resolved Gauss map to the Igusa quartic.  Over a smooth point of \(I\)
the fibre is one reduced point.  The singular locus of \(I\) is the union
of fifteen lines \(L_M\), indexed by matchings
\(M=ij|kl|mn\).  Write a point of \(L_M\), after reordering, as
\((\alpha,\alpha,\beta,\beta,\gamma,\gamma)\) with
\(\alpha+\beta+\gamma=0\).  Away from the intersections of the singular
lines, the resolved fibre is the strict transform of the smooth conic
\[
 (\beta-\gamma)a^2-(\alpha-\gamma)b^2
      +(\alpha-\beta)c^2=0
\]
in the Segre plane
\((a,-a,b,-b,c,-c)\).  This conic passes through the four Segre nodes in
that plane; the fibre in the original rational domain is the conic with
those four base points removed.  At each of the fifteen points where three
singular lines meet, the fibre closure is a union of six projective lines,
and the resolved fibre is the union of their strict transforms.
\end{proposition}

\begin{proof}
Centered squaring is the Gauss map of the Segre cubic.  Segre--Igusa
duality makes it birational, and the fifteen Segre planes are contracted
to the fifteen singular lines of the Igusa quartic
\cite{KondoSegre}.  This proves the smooth-fibre statement and locates the
exceptional fibres.

If the coordinates of \(w\) are paired as above, the corresponding
\(z\)-squares are paired.  Away from an intersection point, the linear and
cubic Segre equations force opposite signs in every pair, so
\(z=(a,-a,b,-b,c,-c)\).  Equality of the centered-square projective
coordinates is exactly the displayed conic.  Its three coefficients are
nonzero precisely away from the three intersection points on \(L_M\).
Putting \(a^2=b^2=c^2\) shows that it contains exactly the four
\(3+3\)-sign nodes of the Segre plane.  These are base points of \(\pi\),
so resolution replaces them by the corresponding points of the strict
transform.
At, for example, \(w=(1,1,1,1,-2,-2)\), the first four \(z_i\) have one
common square and the last two another.  The Segre equations force two
plus and two minus signs among the first four and opposite signs in the
last pair.  Modulo common sign there are six patterns, each with one free
ratio of magnitudes; these are the six lines.  The three pairings of the
first four indices give the three reducible conics through them.
\end{proof}

Thus centered squaring is generically faithful on projective orientation
classes.  Its nonzero information loss is exactly the contraction of the
strict transforms of the fifteen Segre planes.  On a signed affine lift it
also forgets the unavoidable common sign \(z\leftrightarrow-z\).

\section{Rational anomaly synthesis and physical cost}

The anomaly interpretation admits an inverse at the same marked level as
the forward operator theorem.  The inverse is classical invariant theory;
the operator adds the physical normalization and postselection cost.

Write the six path points as \(p_i=(u_i:v_i)\in\mathbf P^1\), and
homogenize a squarefree monomial by
\[
 x_S\longmapsto
 \prod_{i\in S}u_i\prod_{i\notin S}v_i.
\]
The six cubics \(Z_T\) are then multihomogeneous of degree one in every
point.  Their projective class is invariant under the diagonal
\(\operatorname{PGL}_2\)-action.

\begin{theorem}[Golden anomaly inverse]
\label{thm:anomaly-inverse}
Let \(q\in\mathbf P^5(\mathbb Q)\) satisfy
\(\sum_Tq_T=\sum_Tq_T^3=0\).
Then \(q\) has a rational Golden-filter preimage.  On the stable locus its
fibre is one \(\operatorname{PGL}_2\)-orbit, with trivial stabilizer for the
labelled configuration.  Moreover:
\begin{enumerate}[label=\textup{(\roman*)},leftmargin=*]
\item the fifteen vectorlike planes pull back to the fifteen pair-collision
  divisors, and
  \[
    e_5(Z(x))=32\prod_{i<j}(x_j-x_i);
  \]
\item the ten nodes are the closed \(3+3\) collision orbits, while six
  distinct path points map to the smooth nonvectorlike locus; strict
  chirality further removes the six walls \(Z_T=0\);
\item if a normalized rational filter \(y\),
  \(\lVert y\rVert_\infty\leq1\), satisfies \(Z(y)=\lambda q\), then
  \[
    p_T^{(3)}(y)=\frac{\lambda^2q_T^2}{500};
  \]
\item for \(q=(11,-10,-8,5,4,-2)\), the filter
  \[
    (-1,-2/3,-1/3,0,1/3,1)
  \]
  has
  \(p_T^{(3)}=4q_T^2/91125\).  Its primitive integral lift is minimal at
  common-affine height three, but its success is not maximal in the real
  projective fibre.
\end{enumerate}
\end{theorem}

\begin{proof}
The quotient
\((\mathbf P^1)^6/\!/\operatorname{PGL}_2\) is the Segre cubic, and the
colored-triangle cubics are its Joubert coordinates
\cite{HMSVOuter,HMSV}.  We record the inverse because its normalization is
needed below.  Fix
\[
 (p_0,p_1,p_2,p_3,p_4,p_5)=(0,1,\infty,a,b,c)
\]
and, for a Segre point \(z\), put
\[
 \alpha=\frac{z_1-z_3}{2},\quad
 \beta=\frac{z_1+z_5}{2},\quad
 \gamma=-\frac{z_3+z_4}{2},\quad
 \delta=-\frac{z_2+z_3}{2},\quad
 \epsilon=\frac{z_0+z_1}{2}.
\]
On the chart \(\delta\ne\epsilon\), \(\delta c\ne\beta\), set
\[
 c=\frac{\beta-\gamma}{\delta-\epsilon},\qquad
 \lambda=\frac{c(1-c)}{\delta c-\beta},\qquad
 a=c+\lambda\delta,\qquad b=c+\lambda\epsilon.
\]
Then
\[
 ab=c+\lambda\alpha,\qquad
 ac=c+\lambda\beta,\qquad
 bc=c+\lambda\gamma.
\]
Substitution in the six marked cubics gives
\(Z(0,1,\infty,a,b,c)=\lambda z\); the remaining identity is the Segre
equation.  Permuting labels gives the other stable charts.  This also shows
that each stable fibre is one projective-linear orbit.  Three distinct
labelled points kill its stabilizer.  At a node the unique closed orbit has
three points at each of two locations and a one-dimensional torus
stabilizer.

The polynomial \(e_5(Z(x))\) is alternating of degree fifteen.  It is
therefore a scalar multiple of the Vandermonde.  Evaluation at
\((0,1,2,3,4,5)\) gives the scalar 32.  The zero divisor of \(e_5\) on the
Segre cubic is the union of its fifteen planes, each consisting of three
opposite charge pairs.  This proves (i), and GIT stability gives (ii).
Multiplying the fifteen marked matching identities gives the intrinsic
factorization
\[
 \prod_{T<U}(Z_T+Z_U)=-e_5(Z)^3.
\]
Thus each vectorlike plane occurs triply in the product of the fifteen
opposite-pair tests and singly in the reduced \(e_5\)-divisor.
It also explains the Boolean boundary: a sign mask with unequal
multiplicities has at least four coincident points and maps to zero, while
a balanced mask is a \(3+3\) node.

Part (iii) is the Slater determinant identity in
Theorem~\ref{thm:propagation}.  Direct evaluation gives
\(Z=(-4/27)(-q)=(4/27)q\) for the displayed filter, proving its probability
formula.  The height statement is an exact finite census of primitive
integer vectors through height three; the verification supplement records
the search domain and independent replay.
\end{proof}

The last clause hides a useful distinction.  The pole of a projectivity
cuts the real fibre into seven intervals.  On the interval giving the global
maximum, the two endpoint paths remain extreme; after fixing them at
\(\pm1\), the representative has the form
\[
 h_t(u)=\frac{u+t}{tu+1},\qquad -1<t<1.
\]
For the filter in Theorem~\ref{thm:anomaly-inverse}, the common amplitude
gain is
\[
 \kappa(t)=\frac{27(1-t^2)^2}{(3-2t)(9-t^2)}.
\]
Exact Sturm counts on all seven intervals show that the global maximum is
the unique critical point in this chart.  It is the root in \((-1,1)\) of
\[
 t^4-3t^3-24t^2+51t-9=0,
\]
namely \(t=0.194719332570311\ldots\).  The success gain is
\(\kappa(t)^2=1.141022808280276\ldots\).  The maximizing parameter is
irrational, so rational filters have only a supremum.  The rational choice
\(t=1/5\) already gives the filter
\((-1,-7/13,-1/7,1/5,1/2,1)\) and amplitude gain \(486/455\).
Its exact success gain over the integral witness is
\(236196/207025\).
Because the fibre changes only the common amplitude scale, the same
algebraic point simultaneously maximizes every nonzero branch probability
and every fixed positive weighted sum of them.  In particular, a uniformly
chosen protocol succeeds with probability
\(\lambda^2\lVert q\rVert_2^2/3000\).

The same optimization is exact for every smooth real nonvectorlike point.
Choose an affine representative with distinct ordered path coordinates
\(r_1<\cdots<r_6\).  In each of the seven pole chambers, let \(a,b\) be the
two paths whose reciprocal images are extreme.  Up to the fixed initial
amplitude scale, normalization to \([-1,1]\) gives
\[
 \mu_{a,b}(s)=
 \frac{8|s-a|^2|s-b|^2}
 {(b-a)^3\prod_{r_i\notin\{a,b\}}|s-r_i|}.
\]
Its interior critical points satisfy
\[
 \frac2{s-a}+\frac2{s-b}
 -\sum_{r_i\notin\{a,b\}}\frac1{s-r_i}=0.
\]
Clearing denominators gives degree at most four: the degree-five term
cancels as \(2+2-4=0\).  Hence the optimal real filter over any rational
charge point is found by exact comparison of roots of at most seven rational
quartics.

The quotient, its cross-ratio inverse, and rational parametrization of the
six-charge anomaly equations are classical or explicit prior art
\cite{HMSVOuter,CostaDobrescuFox,GripaiosNguyen}.  The assertion here is
the compatibility of that inverse with the marked Golden filter and its
exact Slater cost.  It does not construct gauge fields or a Lagrangian.

\subsection{Two-Abelian anomaly lines}

Let \(q,r\in\mathbb Q^6\) be the charge vectors of two Abelian factors.
The cubic along their projective line is
\[
 \sum_i(q_i+tr_i)^3
 =\sum_iq_i^3+3t\sum_iq_i^2r_i
  +3t^2\sum_iq_ir_i^2+t^3\sum_ir_i^3.
\]
Thus the six individual and mixed anomaly equations say exactly that
\(\mathbf P\langle q,r\rangle\) is a line on the Segre cubic.  We now
transport the classical Fano scheme of these lines through the marked
Golden inverse.

\begin{theorem}[Golden realization of the Fano components]
\label{thm:two-u1-lines}
The Golden control model realizes all 21 components of the Fano scheme of
lines on the Segre cubic as follows.
\begin{enumerate}[label=\textup{(\roman*)},leftmargin=*]
\item The fifteen plane components are the fixed pair-collision divisors
  \(p_i=p_j\).  The frozen duad--syntheme correspondence identifies each
  collision with its three opposite charge pairs.
\item The six remaining components are the six choices of a moving path:
  for fixed \(p_j\), \(j\ne i\),
  \[
    L_i(p_{\widehat\imath})
    =\{Z(p_0,\ldots,p_{i-1},u,p_{i+1},\ldots,p_5):u\in\mathbf P^1\}.
  \]
  Their parameter surfaces are
  \(\overline M_{0,5}\), hence split del Pezzo surfaces of degree five.
\item The five values \(u=p_j\), \(j\ne i\), give the five vectorlike
  crossings of a general chiral line.  Their five charge synthemes form a
  synthematic total and recover the moving-path component.
\item Homogeneous path controls of degree one suffice for every rational
  anomaly line and this degree is minimal for a nonconstant line.  On a
  C715 inverse chart the controls are ratios of linear matching forms.
\item Through a generic smooth nonvectorlike amplitude point pass exactly
  six chiral lines, with directions \(\partial_iZ(x)\), one for each path
  control.  Each direction is itself anomaly-free and satisfies both mixed
  equations with \(Z(x)\).
\end{enumerate}
\end{theorem}

\begin{proof}
Each marked cubic is multihomogeneous of degree one in each path point.
Varying one point therefore traces a line.  Five general fixed points give
a line outside the fifteen planes, and their moduli form \(M_{0,5}\).
The classical Fano classification gives exactly six nonplane components,
each a split degree-five del Pezzo surface \cite{GripaiosNguyenTwo}; their
closures are therefore the six one-moving-path families.  Collision of the
moving point with each fixed point gives the five boundary planes.  The
C715 matching dictionary turns those collisions into the five synthemes
of one total and distinguishes the six families.

A fixed collision maps onto its marked Segre plane, whose lines form the
dual plane.  This gives the other fifteen components.  Finally, restricting
the C715 matching ratios to \(q+tr\) makes every numerator and denominator
linear in \(t\).  A degree-zero control has constant image, proving the
minimality statement.  Multi-affinity gives the six displayed response
directions through \(Z(x)\).  The C715 inverse fibre is one labelled
\(\operatorname{PGL}_2\)-orbit, so a line in the component indexed by \(i\)
through a generic point must fix the same other five path points.  Hence
these six lines are all of them.  Coefficient comparison in the Segre cubic
identity proves the individual and mixed anomaly assertions for each
response direction.
\end{proof}

The source notation for the six nonplane components translates to the C707
path marking as
\[
 (D_0,D_1,D_2,D_3,D_4,D_5)
 \longleftrightarrow
 (D^{\rm path}_4,D^{\rm path}_1,D^{\rm path}_0,
  D^{\rm path}_2,D^{\rm path}_3,D^{\rm path}_5).
\]
For example, moving path 5 has collision total
\[
 \{01|24|35,02|15|34,03|12|45,04|13|25,05|14|23\}.
\]

The control family
\[
 x(s)=\left(0,-\frac12,-1,-\frac37,s,-\frac35\right)
\]
lies in the physical cube for the parameters used below and gives
\[
 Z(x(1/2))=\frac1{140}(59,29,19,-55,-41,-11),
\]
\[
 Z(x(3/4))=\frac3{280}(49,23,17,-45,-35,-9).
\]
Both primitive vectors are strictly chiral, their mixed anomaly sums
vanish, and their ordered charge pairs contain no opposite pair.  A
nonchiral witness is
\[
 y(s)=\left(0,0,-1,-\frac12,\frac12,s\right),
\]
with
\[
 Z(y(2/3))=\frac1{12}(-13,13,-3,3,1,-1),\qquad
 Z(y(3/4))=\frac1{16}(-19,19,-5,5,1,-1).
\]
Here the charge pairs \(01,23,45\) are opposite for both factors.

There is also a direct Majorana form of the two mixed equations.  Put
\(A_T(x)=[D_x,C_T]\).  Theorem~\ref{thm:propagation} gives
\(\Pf A_T(x)=4Z_T(x)\) and
\(\det A_T(x)=16Z_T(x)^2\).  Hence
\[
 \boxed{\sum_T\det A_T(x)\Pf A_T(y)=0},\qquad
 \boxed{\sum_T\Pf A_T(x)\det A_T(y)=0}.
\]
The two left sides are 64 times the mixed charge sums.  Along a one-path
sweep they are the middle coefficients of
\(\sum_T\Pf A_T(x(s))^3=0\).  The line decomposition is classical; the
one-control realization, exact marking, and Pfaffian normalization are the
operator transport.

\section{Verification and scope}

The proofs above separate conceptual arguments from exact finite checks.
The exterior, commutator, golden-compression, frustration, Naimark, recovery,
and anomaly-line mechanisms have human proofs in the text.  Exact arithmetic
fixes normalizations, checks finite orbit tables, proves reducedness of the
fifteen-line spinor base scheme, and certifies the height and real-optimization
clauses of Theorem~\ref{thm:anomaly-inverse}.  The paper-owned verification
supplement records the search domains, hashes, commands, and independent
replays for those claims.  No post-700 theorem in this paper is presently
Lean-verified.

The exceptional-group, Clifford, bad-prime, and lattice constructions are
parents or boundary tests, not additional steps in
Theorem~\ref{thm:propagation}.  Likewise, the finite Majorana family proves
Pfaffian parity walls but does not construct a gauge theory, an interacting
fermionic phase, a protected parity pump, or a topological qubit.  These
boundaries are part of the theorem hierarchy: omitting any optional parent
does not break the operator proof.

\appendix
\section{The collision filtration}
\label{app:collision-filtration}

Work in the translation gauge on the projective axis augmentation space,
and let \(J_Z\) be the \(6\times5\) Jacobian of the six Joubert cubics.
The reduced pair-collision divisor is the union of the fifteen hyperplanes
\(x_i=x_j\).  The scheme \(V(I_4(J_Z))\) has degree \(20\); its radical is
the union of the twenty planes on which a chosen triple of coordinates
coincides.  It is nonreduced, and the annihilator of its nilradical has
radical supported at the ten points indexed by complementary \(3+3\)
partitions.  Finally, the simultaneous-Pfaffian base scheme is the reduced
union of the fifteen lines on which a chosen four coordinates coincide.
Thus their reduced supports form the strict chain
\[
 \{\text{four coincide}\}\subset
 \{\text{three coincide}\}\subset
 \{\text{two coincide}\},
\]
while the Jacobian scheme has an additional lower-dimensional nilpotent
defect at the \(3+3\) points.

The reduced support statements follow from the six-point GIT collision
strata and the radical calculation used in
Proposition~\ref{prop:synchronized-spinors}.  The nilpotent refinement is
currently certified by one exact Singular computation.  It is not used in
the proof of the propagation, rigidity, recovery, or Naimark theorems and is
retained here with that single-CAS trust boundary.

\begin{thebibliography}{9}

\bibitem{KondoSegre}
S.~Kond\=o,
\emph{The Segre cubic and Borcherds products},
arXiv:1110.1126.

\bibitem{HMSV}
B.~Howard, J.~Millson, A.~Snowden, and R.~Vakil,
\emph{The relations among invariants of points on the projective line},
arXiv:0906.2437.

\bibitem{HMSVOuter}
B.~Howard, J.~Millson, A.~Snowden, and R.~Vakil,
\emph{A description of the outer automorphism of \(S_6\), and the
invariants of six points in projective space},
J. Combin. Theory Ser. A 115 (2008), 1296--1303.

\bibitem{CostaDobrescuFox}
D.~B. Costa, B.~A. Dobrescu, and P.~J. Fox,
\emph{General solution to the \(U(1)\) anomaly equations},
Phys. Rev. Lett. 123 (2019), 151601.

\bibitem{GripaiosNguyen}
B.~Gripaios and K.~Le Nguyen Nguyen,
\emph{Anomaly cancellation for a \(U(1)\) factor},
arXiv:2508.11583.

\bibitem{GripaiosNguyenTwo}
B.~Gripaios and K.~Le Nguyen Nguyen,
\emph{Anomaly cancellation for two \(U(1)\) factors},
arXiv:2607.09879.

\bibitem{ChiriviMaffei}
R.~Chiriv\`i and A.~Maffei,
\emph{Pfaffians and shuffling relations for the spin module},
arXiv:1203.2943.

\bibitem{CaiGorenstein}
J.-Y.~Cai and A.~Gorenstein,
\emph{Matchgates revisited},
arXiv:1303.6729.

\bibitem{TschinkelZhang}
Y.~Tschinkel and Z.~Zhang,
\emph{Stable equivariant birationalities of cubic and degree 14 Fano
threefolds},
arXiv:2409.08392.

\end{thebibliography}

\end{document}
